% GENERATED FILE. Edit canonical lessons, then run npm run generate:books.
% canonical-source-hash: fnv1a64:d4c2985dd28eaa43
% canonical-lessons: UR-C338-misal, UR-C338-yani, UR-C338-vazeh, UR-C338-ginana, UR-C338-havala, UR-R338-review-331-338

\chapter{For Example}
\label{ch:ur-for-example}
\hlchaptermodality{338}

\begin{chapteropening}
\textbf{By the end of this chapter:} \emph{I can back up a point with an example.}

Complete the last lesson of chapter 338: I can back up a point with an example.
\end{chapteropening}

\section[\texorpdfstring{\ur{مثال} (misāl)}{مثال (misāl)}]{\ur{مثال} (misāl) — an example}

\label{lesson:UR-C338-misal}



\begin{quote}
Before the new one: say the Urdu for to own up to, then the Urdu for a misunderstanding.
\end{quote}

\subsection*{You'll want to know: \ur{مثال}}
\textbf{\ur{مثال}} — \emph{misāl} — \textquotedblleft{}an example\textquotedblright{}.

\begin{quote}
\ur{ایک} \ur{مثال} \ur{دیجیے۔}
\end{quote}

\begin{quote}
\emph{ek misāl dījiye.}
\end{quote}

\begin{quote}
Give an example, please.
\end{quote}

Feminine.

\begin{cousinweb}
\textbf{\ur{مثال}} (\emph{misāl}) is usually traced, through Persian, to Arabic \emph{mithāl}, \textquotedblleft{}a likeness, a model\textquotedblright{}.
\end{cousinweb}

\begin{grammarlens}[title={What you've built}]
One more word for giving reasons and opinions.
\end{grammarlens}

\subsection*{Your turn}
\begin{itemize}
\raggedright
  \item \emph{Say it:} \emph{misāl}
  \item \emph{Say it:} \emph{misāl}, once more
  \item \emph{Say it:} \emph{ek misāl dījiye.} — Give an example, please.
  \item \emph{Recall:} say the Urdu for to own up to, then the Urdu for a misunderstanding, then say \emph{misāl} again
\end{itemize}

\begin{tcolorbox}[breakable,colback=teal!4,colframe=teal!35!black,title={Before you move on}]
What does \ur{مثال} mean? (\textquotedblleft{}An example\textquotedblright{}.) Say it once more.
\end{tcolorbox}
\section[\texorpdfstring{\ur{یعنی} (yānī)}{یعنی (yānī)}]{\ur{یعنی} (yānī) — that is to say; in other words}

\label{lesson:UR-C338-yani}



\begin{quote}
Before the new one: say the Urdu for a misunderstanding, then the Urdu for an example.
\end{quote}

\subsection*{You'll want to know: \ur{یعنی}}
\textbf{\ur{یعنی}} — \emph{yānī} — \textquotedblleft{}that is to say; in other words\textquotedblright{}.

\begin{quote}
\ur{وہ} \ur{پرسوں} \ur{آئے} \ur{گا،} \ur{یعنی} \ur{اتوار} \ur{کو۔}
\end{quote}

\begin{quote}
\emph{voh parsoṅ āe gā, yānī itvār ko.}
\end{quote}

\begin{quote}
He will come the day after tomorrow, that is, on Sunday.
\end{quote}

\begin{cousinweb}
\textbf{\ur{یعنی}} (\emph{yānī}) is usually traced, through Persian, to Arabic \emph{yaʿnī}, \textquotedblleft{}it means\textquotedblright{}.
\end{cousinweb}

\begin{grammarlens}[title={What you've built}]
One more word for giving reasons and opinions.
\end{grammarlens}

\subsection*{Your turn}
\begin{itemize}
\raggedright
  \item \emph{Say it:} \emph{yānī}
  \item \emph{Say it:} \emph{yānī}, once more
  \item \emph{Say it:} \emph{voh parsoṅ āe gā, yānī itvār ko.} — He will come the day after tomorrow, that is, on Sunday.
  \item \emph{Recall:} say the Urdu for a misunderstanding, then the Urdu for an example, then say \emph{yānī} again
\end{itemize}

\begin{tcolorbox}[breakable,colback=teal!4,colframe=teal!35!black,title={Before you move on}]
What does \ur{یعنی} mean? (\textquotedblleft{}That is to say; in other words\textquotedblright{}.) Say it once more.
\end{tcolorbox}
\section[\texorpdfstring{\ur{واضح} (vāzeh)}{واضح (vāzeh)}]{\ur{واضح} (vāzeh) — clear in meaning; plain}

\label{lesson:UR-C338-vazeh}



\begin{quote}
Before the new one: say the Urdu for an example, then the Urdu for that is to say.
\end{quote}

\subsection*{You'll want to know: \ur{واضح}}
\textbf{\ur{واضح}} — \emph{vāzeh} — \textquotedblleft{}clear in meaning; plain\textquotedblright{}.

\begin{quote}
\ur{آپ} \ur{کا} \ur{جواب} \ur{واضح} \ur{نہیں} \ur{ہے۔}
\end{quote}

\begin{quote}
\emph{āp kā javāb vāzeh nahīṅ hai.}
\end{quote}

\begin{quote}
Your answer is not clear.
\end{quote}

\begin{cousinweb}
\textbf{\ur{واضح}} (\emph{vāzeh}) is usually traced, through Persian, to Arabic \emph{wāḍiḥ}, \textquotedblleft{}clear, evident\textquotedblright{}.
\end{cousinweb}

\begin{grammarlens}[title={What you've built}]
One more word for giving reasons and opinions.
\end{grammarlens}

\subsection*{Your turn}
\begin{itemize}
\raggedright
  \item \emph{Say it:} \emph{vāzeh}
  \item \emph{Say it:} \emph{vāzeh}, once more
  \item \emph{Say it:} \emph{āp kā javāb vāzeh nahīṅ hai.} — Your answer is not clear.
  \item \emph{Recall:} say the Urdu for an example, then the Urdu for that is to say, then say \emph{vāzeh} again
\end{itemize}

\begin{tcolorbox}[breakable,colback=teal!4,colframe=teal!35!black,title={Before you move on}]
What does \ur{واضح} mean? (\textquotedblleft{}Clear in meaning; plain\textquotedblright{}.) Say it once more.
\end{tcolorbox}
\section[\texorpdfstring{\ur{گنانا} (ginānā)}{گنانا (ginānā)}]{\ur{گنانا} (ginānā) — to list; to enumerate (reasons, faults)}

\label{lesson:UR-C338-ginana}



\begin{quote}
Before the new one: say the Urdu for that is to say, then the Urdu for clear in meaning.
\end{quote}

\subsection*{You'll want to know: \ur{گنانا}}
\textbf{\ur{گنانا}} — \emph{ginānā} — \textquotedblleft{}to list; to enumerate (reasons, faults)\textquotedblright{}.

\begin{quote}
\ur{اس} \ur{نے} \ur{تین} \ur{سبب} \ur{گنائے۔}
\end{quote}

\begin{quote}
\emph{us ne tīn sabab gināe.}
\end{quote}

\begin{quote}
He listed three causes.
\end{quote}

Past, with \ur{نے} (he or she …): \ur{اس} \ur{نے} \ur{سبب} \ur{گنائے}, \emph{us ne sabab gināe}.

\begin{cousinweb}
\textbf{\ur{گنانا}} (\emph{ginānā}) is the causative of \textbf{\ur{گننا}}, \textquotedblleft{}to count\textquotedblright{}, which is inherited: CDIAL lists the Hindi cousin \emph{ginnā} under Sanskrit \emph{gaṇáyati}, \textquotedblleft{}counts\textquotedblright{} (CDIAL 3993). To count things out one by one.
\end{cousinweb}

\begin{grammarlens}[title={What you've built}]
One more word for giving reasons and opinions.
\end{grammarlens}

\subsection*{Your turn}
\begin{itemize}
\raggedright
  \item \emph{Say it:} \emph{ginānā}
  \item \emph{Say it:} \emph{ginānā}, once more
  \item \emph{Say it:} \emph{us ne tīn sabab gināe.} — He listed three causes.
  \item \emph{Recall:} say the Urdu for that is to say, then the Urdu for clear in meaning, then say \emph{ginānā} again
\end{itemize}

\begin{tcolorbox}[breakable,colback=teal!4,colframe=teal!35!black,title={Before you move on}]
What does \ur{گنانا} mean? (\textquotedblleft{}To list; to enumerate (reasons, faults)\textquotedblright{}.) Say it once more.
\end{tcolorbox}
\section[\texorpdfstring{\ur{حوالہ} (havālā)}{حوالہ (havālā)}]{\ur{حوالہ} (havālā) — a reference (to a source)}

\label{lesson:UR-C338-havala}



\begin{quote}
Before the new one: say the Urdu for clear in meaning, then the Urdu for to enumerate.
\end{quote}

\subsection*{You'll want to know: \ur{حوالہ}}
\textbf{\ur{حوالہ}} — \emph{havālā} — \textquotedblleft{}a reference (to a source)\textquotedblright{}.

\begin{quote}
\ur{اپنی} \ur{بات} \ur{کا} \ur{حوالہ} \ur{دیجیے۔}
\end{quote}

\begin{quote}
\emph{apnī bāt kā havālā dījiye.}
\end{quote}

\begin{quote}
Give a source for your point, please.
\end{quote}

Masculine.

\begin{cousinweb}
\textbf{\ur{حوالہ}} (\emph{havālā}) is usually traced, through Persian, to Arabic \emph{ḥawāla}, \textquotedblleft{}a handing over\textquotedblright{}; Urdu uses it for pointing a reader on to a source, a sense our sources do not trace step by step.
\end{cousinweb}

\begin{grammarlens}[title={What you've built}]
One more word for giving reasons and opinions.
\end{grammarlens}

\subsection*{Your turn}
\begin{itemize}
\raggedright
  \item \emph{Say it:} \emph{havālā}
  \item \emph{Say it:} \emph{havālā}, once more
  \item \emph{Say it:} \emph{apnī bāt kā havālā dījiye.} — Give a source for your point, please.
  \item \emph{Recall:} say the Urdu for clear in meaning, then the Urdu for to enumerate, then say \emph{havālā} again
\end{itemize}

\begin{tcolorbox}[breakable,colback=teal!4,colframe=teal!35!black,title={Before you move on}]
What does \ur{حوالہ} mean? (\textquotedblleft{}A reference (to a source)\textquotedblright{}.) Say it once more.
\end{tcolorbox}
\section[(dialogue)]{Looking back: from What Caused It to For Example}

\label{lesson:UR-R338-review-331-338}



\begin{quote}
Say the Urdu for a cause and a basis.
\end{quote}

\subsection*{Your turn}
Say these aloud:

\begin{itemize}
\raggedright
  \item a cause, a basis, responsible, to be linked, connected
  \item harm, a means, a reaction, to spoil, to spread
  \item a purpose, for the sake of, a motive, a benefit, to use
  \item to argue, a reasoned argument, to interrupt, an altercation, logic
  \item against, to back, to oppose, in agreement, a side in a dispute
  \item to prove, firm, an accusation, a document, genuine
  \item to confirm, to acknowledge, to give the lie to, to own up to, a misunderstanding
  \item an example, that is to say, clear in meaning, to enumerate, a reference
\end{itemize}

\begin{tcolorbox}[breakable,colback=teal!4,colframe=teal!35!black,title={Before you move on}]
A cause? (\textbf{\ur{سبب}}, \emph{sabab}.) A basis? (\textbf{\ur{بنیاد}}, \emph{bunyād}.) Responsible? (\textbf{\ur{ذمہ} \ur{دار}}, \emph{zimme-dār}.) To be linked? (\textbf{\ur{جڑنا}}, \emph{juṛnā}.) Connected? (\textbf{\ur{وابستہ}}, \emph{vābastā}.) Harm? (\textbf{\ur{نقصان}}, \emph{nuqsān}.) A means? (\textbf{\ur{ذریعہ}}, \emph{zariyā}.) A reaction? (\textbf{\ur{رد} \ur{عمل}}, \emph{radd-e-amal}.) To spoil? (\textbf{\ur{بگاڑنا}}, \emph{bigāṛnā}.) To spread? (\textbf{\ur{پھیلنا}}, \emph{phailnā}.) A purpose? (\textbf{\ur{مقصد}}, \emph{maqsad}.) For the sake of? (\textbf{\ur{خاطر}}, \emph{khātir}.) A motive? (\textbf{\ur{غرض}}, \emph{gharaz}.) A benefit? (\textbf{\ur{فائدہ}}, \emph{fāidā}.) To use? (\textbf{\ur{استعمال} \ur{کرنا}}, \emph{istemāl karnā}.) To argue? (\textbf{\ur{بحث} \ur{کرنا}}, \emph{bahas karnā}.) A reasoned argument? (\textbf{\ur{دلیل}}, \emph{dalīl}.) To interrupt? (\textbf{\ur{ٹوکنا}}, \emph{ṭoknā}.) An altercation? (\textbf{\ur{تکرار}}, \emph{takrār}.) Logic? (\textbf{\ur{منطق}}, \emph{mantiq}.) Against? (\textbf{\ur{خلاف}}, \emph{khilāf}.) To back? (\textbf{\ur{حمایت} \ur{کرنا}}, \emph{himāyat karnā}.) To oppose? (\textbf{\ur{مخالفت} \ur{کرنا}}, \emph{mukhālafat karnā}.) In agreement? (\textbf{\ur{متفق}}, \emph{muttafiq}.) A side in a dispute? (\textbf{\ur{فریق}}, \emph{farīq}.) To prove? (\textbf{\ur{ثابت} \ur{کرنا}}, \emph{sābit karnā}.) Firm? (\textbf{\ur{پختہ}}, \emph{pukhtā}.) An accusation? (\textbf{\ur{الزام}}, \emph{ilzām}.) A document? (\textbf{\ur{دستاویز}}, \emph{dastāvez}.) Genuine? (\textbf{\ur{کھرا}}, \emph{kharā}.) To confirm? (\textbf{\ur{تصدیق} \ur{کرنا}}, \emph{tasdīq karnā}.) To acknowledge? (\textbf{\ur{تسلیم} \ur{کرنا}}, \emph{taslīm karnā}.) To give the lie to? (\textbf{\ur{جھٹلانا}}, \emph{jhuṭlānā}.) To own up to? (\textbf{\ur{قبول} \ur{کرنا}}, \emph{qabūl karnā}.) A misunderstanding? (\textbf{\ur{غلط} \ur{فہمی}}, \emph{ghalat-fahmī}.) An example? (\textbf{\ur{مثال}}, \emph{misāl}.) That is to say? (\textbf{\ur{یعنی}}, \emph{yānī}.) Clear in meaning? (\textbf{\ur{واضح}}, \emph{vāzeh}.) To enumerate? (\textbf{\ur{گنانا}}, \emph{ginānā}.) A reference? (\textbf{\ur{حوالہ}}, \emph{havālā}.)
\end{tcolorbox}
